\scanpage{46}
\mspage{2}
% archive sheet number: 45
\uline{Objets}: Triples $(X,\widetilde\Sigma,\deleted{\unsure{G|S}}\inserted{R})$, où

\margin{$X,\deleted{\unsure{\partial X}},\Sigma,S$}

$X$ surface à bord compacte orientée

$\widetilde\Sigma$ sous-$1$-variété compacte de $X\deleted{\unsure{\setminus\partial X}}$, $\widetilde\Sigma\supset\partial X$

$\deleted{\unsure{G|S}}R\subset\Sigma\deleted{\unsure{\setminus\partial X}}$, partie finie

% slanted
\margin{$\widetilde\Sigma=\partial X\amalg\Sigma$\par On peut considérer que $R$ rigidifie, \unsure{tandis} que $\deleted{R}\inserted{\widetilde\Sigma}$ ``découpe''\unsure{\ldots}}

\uline{flèches}: classes d'isotopie d'homéom. \inserted{conservant l'orientation}

\deleted{Variante} \deleted{\unsure{en}} Les compos. connexes de $\widetilde\Sigma\setminus\partial X\inserted{=\Sigma}$ sont les \uline{cercles de découpage}, \inserted{celles de $\partial X$ sont les cercles de \unsure{raccord|recollement}}; les éléments de $\deleted{\unsure{G|S}}R$ sont les \deleted{\unsure{points marqués}} \uline{sommets de \deleted{découpage}\inserted{rigidification}} \deleted{(ou de rigidification)}, \deleted{les composantes connexes de $\partial X$ sont les ``trous''}. Les composantes connexes de $X$ sont les \deleted{morceaux}\inserted{\uline{composants}} de l'objet envisagé.

Variante
\[
S_*\mathscr{T}\unsure{\mathbb{D}|D}_{\mathrm{top}}^{\pm},\quad\text{avec}
\]
\diagram{\begin{tikzcd}[ampersand replacement=\&]
S_*\mathscr{T}\mathbb{D}_{\mathrm{top}}\arrow[r,phantom,"\subset"] \& S_*\mathscr{T}\mathbb{D}_{\mathrm{top}}^{\pm}\arrow[r] \& {[\pm1]}
\end{tikzcd}}
comme tantôt.

On a des sous-groupoïdes pleins
\diagram{\[
\begin{array}{lll}
R=\varnothing,\ \Sigma=\deleted{\unsure{\partial X}}\varnothing & \mathscr{T}_{\mathrm{top}}\ \text{et}\ \mathscr{T}_{\mathrm{top}}^{\pm} & \\
\deleted{\unsure{G=\widetilde\Sigma=\varnothing}} & & \\
\Sigma=\deleted{\unsure{\partial X}}\varnothing & S_*\mathscr{T}_{\mathrm{top}}\ \text{et}\ S_*\mathscr{T}_{\mathrm{top}}^{\pm} & \text{(\unsure{groupoïdes|groupoïde} de Teichmüller} \\
\text{\inserted{$R$ quelconque}} & & \text{\uline{spéciaux}\inserted{ou \uline{rigidifiés}}} \\
& & \text{(topologiques \unsure{orientés|orienté}))} \\
\deleted{\unsure{G}}R=\varnothing & \mathscr{T}\mathbb{D}_{\mathrm{top}}\ \text{et}\ \mathscr{T}\mathbb{D}_{\mathrm{top}}^{\pm} &
\end{array}
\]}
(Cas pas intéressant je crois, sauf pour $\Sigma=\deleted{\unsure{\partial X}}\varnothing$ également)

% slanted
\margin{$\unsure{(MD)|MD}={}$``modulaire'' ou Mumford-Deligne}

\deleted{\unsure{Dé}} \uline{Détermination} du $\pi_0(S_*\mathscr{T}\mathbb{D}_{\mathrm{top}})\xrightarrow{\sim}\pi_0(S_*\mathscr{T}\mathbb{D}_{\mathrm{top}}^{\pm})$.

On introduit les \uline{graphes} \deleted{\unsure{MD}} bipondérés
